% discussion.tex -- Results 4.5: discussion. 17 Oct (v3_full): rewritten in question order on the freeze results;
% skel_discussion merged (feedforward against feedback with the scored 1-DOF signs; the magnet margin; design rules).
% Every number is a result of Sections 4.1-4.4 or of the literature table's sources below.

\subsection{Discussion}
\label{sec:discussion}

This section discusses the answers in question order, then what the twin can and cannot say, the literature, the
design rules that follow and the limitations.

\paragraph{RQ1: carrying on Betaflight with a measured pull.} The cascade carries the ring with all control on the
ground and the UAVs' firmware unchanged, and what limited the carry on the rig was what
the trackers were told about the rods, together with the rod plan. Measuring the rod's pull at each UAV is the ground-side counterpart of the onboard force
inversion of Sun et al.\ \cite{tudelft}, and the rigid-rod twin shows how far it reaches: it takes up a vertical force
and an error in the typed mass at the UAVs themselves (an exploratory pair separates it), so the planner's integral is
left with the horizontal axes and its vertical channel, which can wind up on an error it cannot see, is now more a
liability than a remedy. At speed the ring cuts the figure-eight's bends by about 3\,cm at 0.3\,m/s; the plan already
carries the cut, and its source is open. On the rig the 3\,cm against 8\,cm came with six settings changed together,
so the measured pull alone is not yet isolated.

\paragraph{RQ2: geometry first, then the pull source.} Several results that first looked like control problems
are geometry: the twin's tilt rise grows as the support margin shrinks and becomes a loss at zero margin, a four-carrier
layout can afford at most two particular leavers, and every rig release that left an even triangle was carried
through. The resize turns this into an advantage for a centralised controller: the planner knows the whole attachment
geometry and solves for the new shares in the first tick after the event, which supports the working hypothesis of
Section~\ref{sec:rq}. What control still sets includes when and how the survivors take up the step, and the one-degree-of-freedom
model frozen before the runs (each survivor's feedforward pull $F = F_{\mathrm{plan}} + k(F_{\mathrm{meas}} -
F_{\mathrm{plan}})$, $k = 0$ on the planner's pull and 1 on the measured pull) has scored its signs on every cell flown so
far: the measured pull gave the smaller tilt rise on an unannounced release in both cells flown (the sign, not yet the
mechanism), detection cut the planner's pull's tilt rise from 28 to 17$^\circ$, detection made no resolvable difference on the measured pull, and the planner's pull tracked a no-event circle
better, narrowly, on the even ring, two runs each. The announced half of the crossover has not flown, so the design rule stands half-tested: the
measured pull in steady flight and on unannounced events, and the plan, ramped, at announced ones only if that half
holds. On the rig the survivors' thrust swung against the step when the plan arrived, in 9 of 9 survivors of the
three unannounced releases, which the twin does not show.

\paragraph{The magnet margin.} The binding hardware constraint is the magnet, not the controller. Every ring lost on
the rig began with a magnet letting go, two of the four after the rods were planned at 60$^\circ$, so the steep plan
reduces the risk without removing it, and the drop rule turned the last loss into a controlled landing. [The hold
against the angle of pull, cold and with the coils on, on a full and a low pack, is still to be measured.]

\paragraph{RQ3: a join is an admission.} Most of the faults that blocked the join in the twin were in the hand-over
between two controllers, not in either controller, which is why the contract makes every observation this work's own;
the rig confirmed it, rejecting its one catch on a proof that assumed the twin's rigid tether. On the rig the attach is
limited earlier, by the magnet reaching the plate. What admission costs the ring waits for the join code tree.

\paragraph{What the resize does not guarantee, and what the twin can say.} Plug-and-play MPC certifies a change before
it is made \cite{riverso2013plug, zeilinger2013plug}; the resize here is shown, not certified: no solve failed at 26
switches, but the only guard against a first solve that goes wrong is the trackers' envelope and the drop rule. On the
releases that predate the bands the twin matched the reactive release and missed what it does not model, and no rig
flight has yet been scored against a frozen band (Section~\ref{sec:ledger}); many apparent
controller faults in simulation were simulator artefacts, from rates differenced on unstamped poses to the engine
ignoring damping on ball joints, and the false touchdown, found the same way, bears on the rig, because a pack or
propeller change can move a UAV's thrust map just as the twin's injected error did.

\paragraph{Against the literature.} Table~\ref{tab:lit_compare} sets this work against the closest works, in the
numbers each reports. On a path the rig's ring stayed within 3\,cm of its reference on average, 3--5\,\% of the
path's half-extent, the same order as Sun et al.'s 4.1\,\% at 1\,m/s and Wang et al.'s 4.7\,\%; but this work flew at
a tenth of their speed, and the twin's
envelope reaches only 0.3\,m/s peak path speed, so the comparison places its accuracy at low
speed only. No work in the table reports a number for the payload's attitude through a departure, or flies a join:
Quan et al.\ release one of four UAVs under a 4\,kg load and report the cable tensions rebalanced after 13\,s, not the
load's attitude. The rig's releases, with tilt rises of 5.3--16.9$^\circ$ over the five carried through, are the only
attitude numbers for a departure in the table; the join has yet to give its number.

\begin{table}[htbp]
\centering
\footnotesize
\setlength{\tabcolsep}{3pt}
\caption{Cooperative transport of a suspended payload: this work against the closest works, as each
reports it. ``--'': not reported.}
\label{tab:lit_compare}
\begin{tabularx}{\textwidth}{@{}
  >{\raggedright\arraybackslash\hsize=0.95\hsize}X
  >{\raggedright\arraybackslash\hsize=0.80\hsize}X
  >{\raggedright\arraybackslash\hsize=0.85\hsize}X
  >{\raggedright\arraybackslash\hsize=1.05\hsize}X
  >{\raggedright\arraybackslash\hsize=1.25\hsize}X
  >{\raggedright\arraybackslash\hsize=0.85\hsize}X
  >{\raggedright\arraybackslash\hsize=1.25\hsize}X@{}}
\toprule
Work & Test; UAVs & Payload; link & Path; speed & Position error\textsuperscript{a} &
Tilt\textsuperscript{b} & UAV removed or added in flight \\
\midrule
Klausen \cite{klausen2020cooperative},\newline 2020 & Real, outdoor; 3 & 1.0, 2.2\,kg point mass; -- &
pilot-set velocity; -- & -- (inter-UAV link error 16--57\,cm) & -- (point mass) & -- \\
\addlinespace
Wang \cite{wangRobustCooperativeTransportation2024},\newline 2024 & Real (mocap); 3 & point mass; cables
0.6\,m (sim) & curve on a 1.3\,m sphere; 3\,m/s (sim) & mean 6.1\,cm in steady state (4.7\,\%) &
-- (point mass) & -- (cable angle 30 to 70$^\circ$ in flight) \\
\addlinespace
Sun \cite{tudelft},\newline 2025 & Real (mocap); 3, some flights 4 & 1.4\,kg; cables 1\,m & figure-eight
$\pm2.5 \times \pm2$\,m; peak 1 to 5\,m/s & RMSE 10.2\,cm (4.1\,\%) at 1\,m/s, 19.7\,cm (7.9\,\%) at
5\,m/s & attitude RMSE 12.9$^\circ$ at 5\,m/s & -- (fixed team) \\
\addlinespace
Quan \cite{quan2025swarm},\newline 2025 (preprint) & Real, outdoor; 4 to 6 & 5\,kg; cables 3 to 4.5\,m &
34\,m delivery; -- & -- (sim: 20--68\,\% of the baselines' error) & -- & one UAV released on command (4 to 3, 4\,kg):
tensions rebalanced after 13\,s, attitude --; 94\,\% success over failures and disconnections; joining not flown \\
\midrule
This work & Real (mocap), twin; 3--4 & 0.86\,kg ring; rigid rods 0.55\,m & circle $r = 0.5$\,m,
0.125\,m/s; figure-eight $\pm1.0 \times \pm0.5$\,m, peak 0.25\,m/s & circle: mean 2.6\,cm (5.2\,\%), max 5.5\,cm;
figure-eight: mean 2.9\,cm (2.9\,\%), max 6.6\,cm & circle 1.7 / 4.0$^\circ$, figure-eight 2.6 / 5.9$^\circ$ (mean /
max) & removed, 4 to 3: tilt rise 5.3--16.9$^\circ$ over five rig releases, loss found in 0.33--0.42\,s; added: [join
tree] \\
\bottomrule
\end{tabularx}

\smallskip
\parbox{\textwidth}{\scriptsize\raggedright
\textsuperscript{a}\,In brackets: the error divided by the path's half-extent along its longer axis
(circle: radius; figure-eight: long semi-amplitude; Wang: sphere radius). This work: mean horizontal
distance from the time-indexed reference (the lag is in it); Sun et al.: 3-D RMSE; Wang et al.:
steady-state mean of the 3-D error. A mean is never larger than the RMSE of the same signal.
\textsuperscript{b}\,Payload tilt. This work: the ring's tilt from level; Sun et al.: axis-angle attitude
error, yaw included. Mocap: motion capture; sim: simulation.
Lee et al.\ \cite{lee2018geometric} and Gabellieri et al.\ \cite{gabellieri2023force} report no position-error number and are left out. This work's tilt is from the day's analysis of each flight; its path errors are on the reviewed window of
Table~\ref{tab:T7}. Not like-for-like: speed (this work 0.125--0.25\,m/s, Sun et al.\ 1--5\,m/s, Wang et al.\ 3\,m/s in
simulation, not restated for the flight); this work hangs the payload on rigid rods, the others on
cables; a point mass (Klausen et al., Wang et al.) has no attitude;
this work's rig entries are one flight each;  Klausen et al.\ and Quan et al.\ fly outdoors on satellite and onboard sensing.}
\end{table}
% ---------------------------------------------------------------------------------------------------
% SOURCES: every number in Table tab:lit_compare. 17 Oct (v3): this work's path errors are now on the T7 window
% (reviewer 12 Oct): circle 2.6 / 5.5 cm (5.2 %), figure-eight 2.9 / 6.6 cm (2.9 %); the lines below keep the 10 Oct registry reading.
% Page references: "p.N" is the printed page; "PDF N" the page index of the file in
% docs/papers/ (multi_drone_control repo). "[computed]" = arithmetic on stated numbers, done here.
%
% Normalisation used for the bracketed % (footnote a): error / half-extent of the path along its longer
% axis. Sun: 2.5 m (x semi-amplitude, Table S1); Wang: 1.3 m (sphere radius, eq. 39); this work: circle
% radius 0.5 m, figure-eight a = 1.0 m (Gerono lemniscate dx = a sin(th), dy = (a/2) sin(2 th), so
% +-1.0 x +-0.5 m; mpc_planner/load_trajectory.py, thesis trajectories.tex).
%
% --- Lee 2018, IEEE TCST 26(1) 255-264, doi 10.1109/TCST.2017.2656060 (file "PDF.js viewer.pdf") ---
%   Simulation only: abstract p.255/PDF 1 "illustrated by a numerical example"; no experiment.
%   n = 3; m0 = 1.5 kg box 1.0 x 0.8 x 0.2 m; m_i = 0.755 kg; l_i = 1 m: p.260/PDF 6, Sec. V.
%   Rigid massless links: abstract ("via rigid links"), Sec. II p.256/PDF 2, Fig. 1 caption.
%   Path x0d = [1.2 sin(0.2 pi t), 4.2 cos(0.1 pi t), -0.5]: p.260/PDF 6. Speed not stated
%   ([computed] peak ~1.5 m/s, not used in the table).
%   Position and attitude error: Fig. 2(a), 2(b) only (Psi0 = 1/2 ||R0 - R0d||^2); no number in the text.
%   No UAV removal/addition; no control rate or computation time.
%
% --- Klausen et al. 2020, IEEE TCST 28(2) 653-660, doi 10.1109/TCST.2018.2876518 (cooperative_control.pdf) ---
%   Real: three hexacopters "each weighing 2.3-kg" (p.653/PDF 1; also m_c = 2.3 kg p.657/PDF 5; p.658
%   says "2.2-kg UAVs": inconsistent in the paper, 2.3 used here). Loads "approximately 1 kg" and
%   "approximately 2.2 kg" (p.657/PDF 5); point-mass load model. Outdoor, GNSS RTK + IMU (p.653, p.657).
%   Formation "equal-sided triangle with sides of 5 m" (p.657/PDF 5); velocity reference set by a
%   pilot (p.657, p.659). Cable length for the flights: not stated (simulation: l_i = 2 m, p.656/PDF 4).
%   Mean link error (inter-UAV distance error, NOT payload error), Table I p.659/PDF 7: light load
%   0.38 / 0.37 / 0.16 m, heavy load 0.57 / 0.34 / 0.46 m -> "16-57 cm" in the table.
%   Simulation (p.656/PDF 4): N = 3, UAV 2 kg, load 3 kg, wires 2 m, 1 m/s, wind 4 m/s; link error
%   plotted only (Fig. 5). Wind on the test day "estimated 4-6 m/s" (p.658/PDF 6).
%   Decentralised passivity-based formation control (eqs. 16-23; Theorem 1 p.655). No rate, no
%   computation time. No UAV removal/addition.
%
% --- Gabellieri et al. 2023, IROS, pp. 6920-6926, doi 10.1109/IROS55552.2023.10342240 (force_based.pdf) ---
%   Real (Sec. VI, p.6925/PDF 6): two quadrotors 1.03 kg, cables 1 m, 2 m bar 0.2 kg; motion capture
%   120 Hz fused with IMU 1 kHz; pose regulation to p_L = [0.5 0 1.5] m, psi 0.2 rad, theta -0.12 rad,
%   t_L in {1, 2, 3} N (three runs). Attitude error: Fig. 5 box plots only (mean e_RL over the last
%   20 s); text qualitative (p.6926/PDF 7). Position error: not reported.
%   Simulation (Sec. V, p.6924/PDF 5): Gazebo, 0.5 kg 1 m bar, 27 runs; only stated error number is
%   "100 deg in Scenario2" (equilibrium attitude error with a large force bias): not a tracking result.
%   Admittance leader-follower, "does not require explicit communication" (Remark 1, p.6921-6922).
%   No rate, no computation time. No UAV removal/addition.
%
% --- Wang et al. 2024, IEEE T-ITS 25(9) 11833-11843 (robust_cooperative.pdf) ---
%   Simulation (Sec. VI, p.11840-11841/PDF 8-9): n = 3; m_L = 0.06 kg (point mass); m_i = 0.21 kg;
%   cables 0.6 m (rigid, fixed length); Clelia curve eq. (39), 1.3 m sphere, constant speed 3 m/s.
%   Converges "within approximately 5 s"; "converges to zero by the fourth second" vs 19 s for [9].
%   Real (Sec. VII, p.11841-11842/PDF 9-10): VICON 100 Hz, three quadrotors (Fig. 1, Fig. 9 legend);
%   tracks eq. (39) (speed not restated; ~1 m/s read off Fig. 8 by the extraction, not stated);
%   "the mean position error is only 6.1 cm in steady-state" (p.11842/PDF 10, Sec. VII-B; Fig. 8
%   inset 0.061 m). [computed] 0.061 / 1.3 = 4.7 %.
%   Masses and cable length for the flights: not restated (sim values shown in the table, marked).
%   Cable-angle reconfiguration theta_d 30 -> 70 -> 30 deg, eq. (40), window flight 0.6 m/s for 5 s
%   (p.11842/PDF 10, Fig. 9); no position-error number for it. No UAV removal/addition.
%   Backstepping (Fig. 5 p.11836/PDF 4); no rate, no computation time.
%
% --- Sun et al. 2025 (bib key tudelft), arXiv 2501.18802v2 (sun2025_...pdf); PDF page index ---
%   Three quadrotors, 1.4 kg payload, quadrotor 0.6 kg, cables 1 m: PDF 5 (Results, Experimental
%   Setup); Fig. S3 PDF 41. Four quadrotors in real flight: PDF 18 ("scaled up ... to four units"),
%   Movie S3 PDF 46; nine in a simulated setpoint task (PDF 18, Fig. 7C).
%   Paths: Table S1 PDF 45: Slow 2.5 cos(0.25t), 2 sin(0.5t); Fast 2.5 cos(t), 2 sin(2t); z 1.0.
%   Table 1 PDF 6: vmax Slow 1, Medium 2, Medium Plus 2, Fast 5 m/s. Real position RMSE: Slow 0.102,
%   Medium 0.093, Medium Plus 0.117, Fast 0.197 m. [computed] 0.102 / 2.5 = 4.1 %, 0.197 / 2.5 = 7.9 %.
%   Fast real attitude RMSE 12.9 deg, PDF 8; axis-angle, Fig. 2 caption PDF 7; yaw reference
%   0.25 rad/s, PDF 6.
%   Planner: central OCP at 10 Hz on a laptop (PDF 6; Fig. 8 PDF 21), horizon 2 s (PDF 5), N = 20
%   (PDF 23), acados SQP-RTI (PDF 23); onboard tracker 300 Hz with INDI (PDF 26).
%   "The average CPU time consumed by the planner was 15.3 ms" (3 quadrotors, Fast, real; PDF 18),
%   Intel i7-13700H; CPU time "grew exponentially" with team size, nine units within 10 Hz (PDF 18).
%   Wind (PDF 14), not in the table: RMSE 0.048 vs 0.055 m (3 UAVs), 0.048 vs 0.070 m (4 UAVs).
%   UAV failure, removal or addition: not tested anywhere in the 46 pages.
%
% --- Quan et al. 2025, arXiv 2509.03563v1, preprint (quan2025_...pdf; the full 33-page copy since 11 Oct) ---
%   Unplug-and-play experiment (PDF 15-17, Fig. 6E-H): four robots, 4 kg, each with a cable release; Robot 2 releases
%   at 40 s; the other three rebalance to a near-equilateral layout, tensions balanced after 13 s; Robot 1's velocity
%   plotted. NO load attitude or swing number anywhere in the paper. Joining: prose only (PDF 7-8 and 19, 'connected
%   to the payload and its neighbors'); varying team sizes in simulation are referred to S5 (supplementary).
%   Teams of four, five and six (PDF 8, Results); 5 kg payload with five robots (PDF 8); 94 % success
%   "under single-robot failure, disconnection events" with six robots (abstract, PDF 1-2; trial
%   count not given). Cable labels 3, 3.5, 4, 4.5 m in the cable-length-uncertainty panel (Fig. 2,
%   PDF 9). Delivery "34 m away", load at "15 m above the ground" (PDF 8). Winds up to Beaufort 4
%   (abstract). Removal shown 4 -> 3 (Fig. 2 Unplug-and-Play, PDF 9); joining described only as a
%   concept (PDF 8). Simulation tracking errors "20%, 68%, 55.5%, and 21.9% of existing approaches"
%   (abstract, PDF 1): relative only. Decentralised force model, no explicit
%   communication (abstract); onboard LiDAR, camera, DGPS (Fig. 2, PDF 9). No absolute error, tilt,
%   transient, rate or computation time in PDF 1-9.
%
% --- This work: results/registry.csv (multi_drone_control repo) ---
%   Payload 0.86 kg ring, rods 0.55 m, UAVs 0.55 kg: supplied with the task (the rig and rig-twin
%   constants; not a registry measurement).
%   Circle: RIG-1007-r206e60 (rig, 4 UAVs, carry mode, even ring, orbit r 0.5 at 0.125 m/s, from +4 s
%   to LAND): xy path err mean 2.91 cm, max 5.7 cm, lag 0.20 s, ring tilt mean 1.7 / max 4.0 deg,
%   ring z sd 0.5 cm, z err -0.4 cm. [computed] 2.91 / 50 = 5.8 %.
%   Figure-eight: RIG-1007-r2008e60 (rig, 4 UAVs, mpc, fig_8 a 1.0 at 0.25 m/s): mean 3.0 cm, max
%   6.8 cm, lag 0.21 s, tilt mean 2.6 / max 5.9 deg, ring z err -0.7 cm. [computed] 3.0 / 100 = 3.0 %.
%   Path error definition: results/rig/2026-10-07/figures/compare_paths.py: horizontal distance of the
%   ring (mocap) from the planner's time-indexed payload reference, mean from +4 s on the path to LAND;
%   lag = the time shift that best aligns the two paths. One flight each.
%   Detach, rig, 4 -> 3, detach mode, drop_on_loss armed, one flight each:
%     RIG-1008-r302_3 announced, hover: tilt peak 15.4 deg, ring z drop -6.9 cm, xy swing 11.9 cm,
%       no drop, landed.
%     RIG-1008-r303_4 unannounced, hover: detected 0.42 s after release, tilt peak 6.5 deg, no drop.
%     RIG-1008-r304_3 unannounced, 150 deg gap: detected 0.33 s, tilt peak 17.8 deg, z drop -5.3 cm,
%       xy swing 16.1 cm, no drop.
%     RIG-1008-r305 unannounced, during the r 0.5 m / 0.125 m/s circle: detected 0.34 s, tilt peak
%       8.1 deg, 3-D RMSE 6.8 cm (a constant ~6 cm offset: detach runs fly int_mode reference), no drop.
%     -> tilt peak 6.5-17.8 deg (4 flights), height dip 5.3-6.9 cm (2 flights that report it),
%        detection 0.33-0.42 s (3 unannounced flights).
%     Twin bands for the same configs: R1207/R1212 7.4/8.6 deg, R1206/R1211 15.5/18.2 deg,
%     R1208/R1213 7.3/6.0 deg.
%     CORRECTION 9 Oct on all four rows: the freed UAV's touchdown detector stalled at 0.17-0.21 m and
%     it was dropped ~0.3 m by the fleet disarm (card 2026-10-09_departed_ground_idle). Not a payload
%     number, so not in the table; it qualifies "landed" for the freed UAV.
%   Attach (no rig attach flown yet): twin R1266 (twin_r8b, exploratory, one run, hover, QUAD4 welded
%     on plate 9, OCP resized 3 -> 4, 3.3 s on four then LAND): ring tilt peak weld..LAND 2.68 deg,
%     ring dip 0.001 m.
%   Planner solve time on the rig (supplied with the task; from the 7-8 Oct rig planner logs,
%     results/rig/2026-10-0[78]/*_logs, tools/thesis_rig_figures.py fig_solve, docs/figures/F_solve_time;
%     not a registry row): median 31 ms (3 UAVs), 66 ms (4 UAVs, detach mode), 51 ms (4 UAVs, mpc
%     mode); no solve over the 100 ms budget. Note: trajectories.tex (Table tab:solve_times) still
%     says 24-30 ms from the 1 Oct flights and that four UAVs were not yet timed on the rig: update it.
%   10 Hz planner, 2 s horizon, per-UAV tracking MPC: the method chapter.
%


\paragraph{Design rules.} Choose a layout whose survivors keep a positive support margin for every planned leaver,
and five carriers if any single loss must be survivable; plan the rods steep and feed each tracker the pull it
measures; keep a rule that drops the ring and lands every UAV in every mode; admit a newcomer only on a geometric
confirmation this controller observes itself, from a small tension; and re-arm a landing's touchdown test when the
load comes off.

\paragraph{Limitations.} The work is at lab scale (a 0.86\,kg payload, three to four UAVs, a 6\,$\times$\,6\,m room,
peak path speeds up to 0.3\,m/s), with motion capture as the only sensing; the twin has only its bridge's modelled latency, not the rig's. The
rods are rigid, so slack cables are outside the model, and the twin's links
are rigid too, the newcomer's tether included, where the rig's go slack. The rig results are one flight per condition,
on code trees that moved between rig days (the 8 October releases flew the reference integral, the current default is
the model integral), and the twin predicted none of them under a frozen band yet. The releases are commanded magnet
releases; a carrier that dies and is left hanging on the ring is out of scope and left to the drop rule.  The join depends on the collaborator's
controller and has not completed on the rig or been scored on the current confirmation. The touchdown re-arm is not yet
on the rig, and a UAV whose throttle offset is 0.03 below its map's still disarms about 16\,cm up in the twin. The vertical integral
can wind up without effect, and the timing fix that removes the corner cutting failed on the spin, so the cut is
reported, not removed.
